% !TEX root = base-r.tex

{\setbeamercolor{block body}{fg = black, bg = white}
\begin{block}{Lists}
  \centering
  \inlc{l <- list(x = 1:5, y = c('a', 'b'))}\\{\small A list is a collection of elements which can be of different types.}
  
  \smallskip
  {\color{gray}\hrule}
  
  \smallskip
  \begin{tabular}{*{4}{>{\centering\arraybackslash}m{0.215\linewidth}}}
    \inlc{l[[2]]} & \inlc{l[1]} & \inlc{l\$x} & \inlc{l['y']}\\
    Second element of l. & New list with only the first element. & Element named x. & New list with only element named y.
  \end{tabular}
\end{block}
}

\begin{block}{Matrices}
  \begin{column}[t]{0.5\linewidth}
    \begin{tabular}{>{\centering\arraybackslash} m{0.65cm} c @{-~~} l}
      \begin{tikzpicture}[scale = 0.3]
        \foreach \x in {0, 1, 2}
          \foreach \y in {0, 2}
            \filldraw[lightgray!65] (\x, \y) rectangle ({\x + 0.8}, {\y + 0.8}); 
        \foreach \x in {0, 1, 2}
          \filldraw[yellow!75!orange] (\x, 1) rectangle ({\x + 0.8}, {1 + 0.8}); 
      \end{tikzpicture}
      & \inlc{m[2,  ]} & Select a row\\
      \begin{tikzpicture}[scale = 0.3]
        \foreach \y in {0, 1, 2}
          \foreach \x in {1, 2}
            \filldraw[lightgray!65] (\x, \y) rectangle ({\x + 0.8}, {\y + 0.8}); 
        \foreach \y in {0, 1, 2}
          \filldraw[yellow!75!orange] (0, \y) rectangle ({0 + 0.8}, {\y + 0.8}); 
      \end{tikzpicture}
      & \inlc{m[ ,1]} & Select a column\\
      \begin{tikzpicture}[scale = 0.3]
        \foreach \x in {0, 1, 2}
          \foreach \y in {0, 1, 2}
            \filldraw[lightgray!65] (\x, \y) rectangle ({\x + 0.8}, {\y + 0.8}); 
        \filldraw[yellow!75!orange] (2, 1) rectangle ({2 + 0.8}, {1 + 0.8}); 
      \end{tikzpicture}
      & \inlc{m[2,3]} & Select an element
    \end{tabular}
  \end{column}          
  {\color{gray}\vrule}          
  \begin{column}[t]{0.35\linewidth}
    \vspace{-10ex}
    \flushright
    \inl{t(m)}\\ Transpose
    
    \inl{m \%*\% n}\\{ Matrix Multiplication}
    
    \inl{solve(m,n)}\\{ Find \inl{x} in \inl{m*x = n}}
  \end{column}
\end{block}

{\setbeamercolor{block body}{fg = black, bg = white}
\begin{seeblock}{Data Frames}{dplyr}
  {\fontsize{9pt}{11pt}\selectfont\inlc{df <- data.frame(x = 1:3, y = c('a', 'b', 'c'))}\\A special case of a list where all elements are the same length.}
  
  \begin{column}{0.35\linewidth}
    
    \renewcommand{\arraystretch}{1.75}
    \color{secondary}
    \begin{tabular}{| >{\centering\arraybackslash\color{black}} m{0.75cm} | >{\centering\arraybackslash\color{black}} m{0.75cm} |}
      \hline
      \rowcolor{codebg}\color{white} x & \color{white}y\\\hline
      \rowcolor{white} 1 & a\\\hline
      \rowcolor{white} 2 & b\\\hline
      \rowcolor{white} 3 & c\\\hline
    \end{tabular}
  \end{column}
%          
  \begin{column}{0.65\linewidth}
    \vspace{-3ex}
    \begin{center}\textbf{List subsetting}\end{center}
    \begin{tabular}{c >{\centering\arraybackslash} m{0.65cm} c m{0.65cm}}
      \inlc{df\$x} & 
      \begin{tikzpicture}[xscale = 0.4, yscale = 0.3]
        \foreach \y in {0, 1, 2}
          \filldraw[yellow!75!orange] (0, \y) rectangle ({0 + 0.8}, {\y + 0.8});
        \filldraw[yellow!70!black] (0, 3) rectangle ({0 + 0.8}, {3 + 0.8}); 
        \foreach \y in {0, 1, 2}
          \filldraw[white] (1, \y) rectangle ({1 + 0.8}, {\y + 0.8});
        \filldraw[codebg] (1, 3) rectangle ({1 + 0.8}, {3 + 0.8}); 
      \end{tikzpicture} &
      \inlc{df[[2]]} &
      \begin{tikzpicture}[xscale = 0.4, yscale = 0.3]
        \foreach \y in {0, 1, 2}
          \filldraw[yellow!75!orange] (1, \y) rectangle ({1 + 0.8}, {\y + 0.8});
        \filldraw[yellow!70!black] (1, 3) rectangle ({1 + 0.8}, {3 + 0.8}); 
        \foreach \y in {0, 1, 2}
          \filldraw[white] (0, \y) rectangle ({0 + 0.8}, {\y + 0.8});
        \filldraw[codebg] (0, 3) rectangle ({0 + 0.8}, {3 + 0.8}); 
      \end{tikzpicture}
    \end{tabular}
    \vspace{-1.5ex}
    \begin{center}
      \begin{tikzpicture}[remember picture, overlay]
        \draw[darkgray, dotted, very thick, rounded corners] (-3, 0) rectangle (3, -1.5);
      \end{tikzpicture}
      
      \textit{Understanding a data frame}
      
      \smallskip
      \begin{tabular}{c l}
        \inl{view(df)} & See the full data frame.\\
        \inl{head(df)} & See the first 6 rows.
      \end{tabular}
    \end{center}
  \end{column}
  
  \bigskip
  \textbf{Matrix subsetting}
  
  \medskip
  \begin{column}{0.33\linewidth}
    \begin{tabular}{c >{\centering\arraybackslash} m{0.65cm}}
      \inlc{df[ ,2]} &
      \begin{tikzpicture}[xscale = 0.4, yscale = 0.3]
        \foreach \y in {0, 1, 2}
          \filldraw[yellow!70!orange] (1, \y) rectangle ({1 + 0.8}, {\y + 0.8});
        \filldraw[yellow!75!black] (1, 3) rectangle ({1 + 0.8}, {3 + 0.8}); 
        \foreach \y in {0, 1, 2}
          \filldraw[white] (0, \y) rectangle ({0 + 0.8}, {\y + 0.8});
        \filldraw[codebg] (0, 3) rectangle ({0 + 0.8}, {3 + 0.8}); 
      \end{tikzpicture}\\[0.75ex]
      \inlc{df[2,  ]} &
      \begin{tikzpicture}[xscale = 0.4, yscale = 0.3]
        \foreach \x in {0, 1} {
          \filldraw[codebg] (\x, 3) rectangle ({\x + 0.8}, {3 + 0.8}); 
          \filldraw[white] (\x, 0) rectangle ({\x + 0.8}, {0 + 0.8}); 
          \filldraw[white] (\x, 2) rectangle ({\x + 0.8}, {2 + 0.8}); 
          \filldraw[yellow!75!orange] (\x, 1) rectangle ({\x + 0.8}, {1 + 0.8}); 
        }
      \end{tikzpicture}\\[0.75ex]
      \inlc{df[2,2]} &
      \begin{tikzpicture}[xscale = 0.4, yscale = 0.3]
        \foreach \x in {0, 1} {
          \filldraw[codebg] (\x, 3) rectangle ({\x + 0.8}, {3 + 0.8}); 
          \foreach \y in {0, 1, 2}
            \filldraw[white] (\x, \y) rectangle ({\x + 0.8}, {\y + 0.8}); 
        }
        \filldraw[yellow!75!orange] (1, 1) rectangle ({1 + 0.8}, {1 + 0.8}); 
      \end{tikzpicture}\\
    \end{tabular}
  \end{column}
  {\color{gray}{\vrule}}
  \begin{column}{0.24\linewidth}
    \vspace{-13ex}
    
    \inl{nrow(df)}\\Number of rows.
    
    \inl{ncol(df)}\\Number of columns.
    
    \inl{dim(df)}\\Number of columns and rows.
  \end{column}
  {\color{gray}{\vrule}}
  \begin{column}{0.4\linewidth}
    \vspace{-12ex}
    
    \inl{cbind} - Bind columns.
    
    \smallskip
    {\centering
    \begin{tikzpicture}[xscale = 0.4, yscale = 0.3]
      \foreach \x in {0, 1} {
        \filldraw[codebg] (\x, 3) rectangle ({\x + 0.8}, {3 + 0.8}); 
        \foreach \y in {0, 1, 2}
          \filldraw[white] (\x, \y) rectangle ({\x + 0.8}, {\y + 0.8}); 
      }
      
      \foreach \y in {0, 1, 2}
        \filldraw[red] (2.5, \y) rectangle ({2.5 + 0.8}, {\y + 0.8}); 
      \filldraw[red!70!black] (2.5, 3) rectangle ({2.5 + 0.8}, {3 + 0.8}); 
      
      \draw[->, > = stealth', ultra thick, main] (3.75, 1.5) -- (5.35, 1.5);
      
      \foreach \x in {6, 7, 8} {
        \foreach \y in {0, 1, 2}
          \filldraw[green] (\x, \y) rectangle ({\x + 0.8}, {\y + 0.8}); 
        \filldraw[green!70!black] (\x, 3) rectangle ({\x + 0.8}, {3 + 0.8}); 
      }
    \end{tikzpicture}}
    
    \smallskip
    \inl{rbind} - Bind rows.
    
    \smallskip
    {\centering
    \begin{tikzpicture}[xscale = 0.4, yscale = 0.3]
      \foreach \x in {0, 1} {
        \filldraw[codebg] (\x, 3) rectangle ({\x + 0.8}, {3 + 0.8}); 
        \foreach \y in {0, 1, 2}
          \filldraw[white] (\x, \y) rectangle ({\x + 0.8}, {\y + 0.8}); 
      }
      
      \foreach \x in {2.5, 3.5} {
        \filldraw[red] (\x, 1) rectangle ({\x + 0.8}, {1 + 0.8}); 
        \filldraw[red!70!black] (\x, 2) rectangle ({\x + 0.8}, {2 + 0.8}); 
      }
      
      \draw[->, > = stealth', ultra thick, main] (4.75, 1.5) -- (6.35, 1.5);
      
      \foreach \x in {7, 8} {
        \foreach \y in {-0.5, 0.5, 1.5, 2.5}
          \filldraw[green] (\x, \y) rectangle ({\x + 0.8}, {\y + 0.8}); 
        \filldraw[green!70!black] (\x, 3.5) rectangle ({\x + 0.8}, {3.5 + 0.8}); 
      }
      
    \end{tikzpicture}}
        
  \end{column}
\end{seeblock}
}